\documentclass[10pt,a4paper]{amsart}
\usepackage{amsmath,amssymb,amsthm,graphicx}
\usepackage{verbatim}
\newenvironment{reference}{\comment}{\endcomment}
%\newenvironment{reference}{}{}
\newenvironment{slogan}{\comment}{\endcomment}
\newenvironment{history}{\comment}{\endcomment}

% For commutative diagrams we use Xy-pic
\usepackage[all]{xy}

% We use 2cell for 2-commutative diagrams.
\xyoption{2cell}
% Theorem environments.
%
\theoremstyle{plain}
\newtheorem{theorem}[subsection]{Theorem}
\newtheorem{proposition}[subsection]{Proposition}
\newtheorem{lemma}[subsection]{Lemma}

\theoremstyle{definition}
\newtheorem{definition}[subsection]{Definition}
\newtheorem{example}[subsection]{Example}
\newtheorem{exercise}[subsection]{Exercise}
\newtheorem{situation}[subsection]{Situation}

\theoremstyle{remark}
\newtheorem{remark}[subsection]{Remark}
\newtheorem{remarks}[subsection]{Remarks}

\numberwithin{equation}{subsection}

% Macros
%
\def\lim{\mathop{\mathrm{lim}}\nolimits}
\def\colim{\mathop{\mathrm{colim}}\nolimits}
\def\Spec{\mathop{\mathrm{Spec}}}
\def\Hom{\mathop{\mathrm{Hom}}\nolimits}
\def\Ext{\mathop{\mathrm{Ext}}\nolimits}
\def\SheafHom{\mathop{\mathcal{H}\!\mathit{om}}\nolimits}
\def\SheafExt{\mathop{\mathcal{E}\!\mathit{xt}}\nolimits}
\def\Sch{\mathit{Sch}}
\def\Mor{\mathop{\mathrm{Mor}}\nolimits}
\def\Ob{\mathop{\mathrm{Ob}}\nolimits}
\def\Sh{\mathop{\mathit{Sh}}\nolimits}
\def\NL{\mathop{N\!L}\nolimits}
\def\CH{\mathop{\mathrm{CH}}\nolimits}
\def\proetale{{pro\text{-}\acute{e}tale}}
\def\etale{{\acute{e}tale}}
\def\QCoh{\mathit{QCoh}}
\def\Ker{\mathop{\mathrm{Ker}}}
\def\Im{\mathop{\mathrm{Im}}}
\def\Coker{\mathop{\mathrm{Coker}}}
\def\Coim{\mathop{\mathrm{Coim}}}

% Boxtimes
%
\DeclareMathSymbol{\boxtimes}{\mathbin}{AMSa}{"02}

%
% Macros for moduli stacks/spaces
%
\def\QCohstack{\mathcal{QC}\!\mathit{oh}}
\def\Cohstack{\mathcal{C}\!\mathit{oh}}
\def\Spacesstack{\mathcal{S}\!\mathit{paces}}
\def\Quotfunctor{\mathrm{Quot}}
\def\Hilbfunctor{\mathrm{Hilb}}
\def\Curvesstack{\mathcal{C}\!\mathit{urves}}
\def\Polarizedstack{\mathcal{P}\!\mathit{olarized}}
\def\Complexesstack{\mathcal{C}\!\mathit{omplexes}}
% \Pic is the operator that assigns to X its picard group, usage \Pic(X)
% \Picardstack_{X/B} denotes the Picard stack of X over B
% \Picardfunctor_{X/B} denotes the Picard functor of X over B
\def\Pic{\mathop{\mathrm{Pic}}\nolimits}
\def\Picardstack{\mathcal{P}\!\mathit{ic}}
\def\Picardfunctor{\mathrm{Pic}}
\def\Deformationcategory{\mathcal{D}\!\mathit{ef}}
\usepackage{cleveref}
\crefname{originalreference}{original reference}{original references}
\Crefname{originalreference}{Original reference}{Original references}
\makeatletter
\newlabel{stacks-sheaves-lemma-pushforward-restriction}{{075G}{source}{}{sourcepointer-1}{}}
\newlabel{stacks-sheaves-lemma-pushforward-restriction@cref}{{[originalreference][0][]075G}{[1][0][]source}{}{}{}}
\makeatother
\crefname{originallemma}{original lemma}{original lemmas}
\Crefname{originallemma}{Original lemma}{Original lemmas}
\makeatletter
\newlabel{lemma-parasitic}{{0774}{source}{}{sourcepointer-2}{}}
\newlabel{lemma-parasitic@cref}{{[originallemma][0][]0774}{[1][0][]source}{}{}{}}
\makeatother
\crefname{originalreference}{original reference}{original references}
\Crefname{originalreference}{Original reference}{Original references}
\makeatletter
\newlabel{stacks-sheaves-proposition-smooth-covering-compute-cohomology}{{06XJ}{source}{}{sourcepointer-3}{}}
\newlabel{stacks-sheaves-proposition-smooth-covering-compute-cohomology@cref}{{[originalreference][0][]06XJ}{[1][0][]source}{}{}{}}
\makeatother
\crefname{originallemma}{original lemma}{original lemmas}
\Crefname{originallemma}{Original lemma}{Original lemmas}
\makeatletter
\newlabel{lemma-general-pushforward}{{076Z}{source}{}{sourcepointer-4}{}}
\newlabel{lemma-general-pushforward@cref}{{[originallemma][0][]076Z}{[1][0][]source}{}{}{}}
\makeatother
\crefname{originalreference}{original reference}{original references}
\Crefname{originalreference}{Original reference}{Original references}
\makeatletter
\newlabel{descent-lemma-cohomology-parasitic}{{0755}{source}{}{sourcepointer-5}{}}
\newlabel{descent-lemma-cohomology-parasitic@cref}{{[originalreference][0][]0755}{[1][0][]source}{}{}{}}
\makeatother
\title{Cohomology of parasitic modules vanishes and pushforwards of parasitic modules are parasitic}
\author{The Stacks Project Authors}
\date{stacks-project:stacks-cohomology.tex:local-snapshot-20261012; GNU FDL 1.2 or later}
\begin{document}
\maketitle
\noindent\textit{Editorial note.} All delivered mathematical text is the author's own, copied verbatim from the pinned Stacks Project chapter file stacks-cohomology.tex; no statement or proof step was rewritten, reordered or removed.\par\medskip
\noindent\textit{Editorial note.} The counted span cites the lemmas recorded in assembly.external\_references with canonical tags as navigation pointers.\par\medskip
\noindent\textit{Allowed standard prerequisites (editorial).}\begin{itemize}
\item Pushforward and restriction formula. Sheaves on Stacks, Lemma stacks-sheaves-lemma-pushforward-restriction. Describes R\textasciicircum{}i f\_* F as a sheafification of stalkwise cohomology.
\item Smooth covering computes cohomology. Sheaves on Stacks, Proposition stacks-sheaves-proposition-smooth-covering-compute-cohomology. Spectral sequence used twice to reduce the stack to a scheme.
\item Cohomology of parasitic modules on schemes. Descent, Lemma descent-lemma-cohomology-parasitic. Covers the final scheme-representable case.
\end{itemize}\medskip
\ifdefined\hypertarget\hypertarget{sourcepointer-1}{}\fi
\noindent\textit{Source reference.} Original reference 075G: sources/stacks-chow/tags.txt, source lines 9130--9130. Cross-chapter standard prerequisite cited by the author inside the counted span; it is used as a pointer and its own proof is not part of this item.\par\medskip
\ifdefined\hypertarget\hypertarget{sourcepointer-2}{}\fi
\noindent\textit{Source reference.} Original lemma 0774: sources/stacks-chow/tags.txt, source lines 9189--9189. Cross-chapter standard prerequisite cited by the author inside the counted span; it is used as a pointer and its own proof is not part of this item.\par\medskip
\ifdefined\hypertarget\hypertarget{sourcepointer-3}{}\fi
\noindent\textit{Source reference.} Original reference 06XJ: sources/stacks-chow/tags.txt, source lines 8853--8853. Cross-chapter standard prerequisite cited by the author inside the counted span; it is used as a pointer and its own proof is not part of this item.\par\medskip
\ifdefined\hypertarget\hypertarget{sourcepointer-4}{}\fi
\noindent\textit{Source reference.} Original lemma 076Z: sources/stacks-chow/tags.txt, source lines 9184--9184. Cross-chapter standard prerequisite cited by the author inside the counted span; it is used as a pointer and its own proof is not part of this item.\par\medskip
\ifdefined\hypertarget\hypertarget{sourcepointer-5}{}\fi
\noindent\textit{Source reference.} Original reference 0755: sources/stacks-chow/tags.txt, source lines 9119--9119. Cross-chapter standard prerequisite cited by the author inside the counted span; it is used as a pointer and its own proof is not part of this item.\par\medskip
\begin{lemma}
\label{lemma-pushforward-parasitic}
Let $\tau \in \{\etale, fppf\}$.
Let $\mathcal{X}$ be an algebraic stack.
Let $\mathcal{F}$ be a parasitic object of
$\textit{Mod}(\mathcal{X}_\tau, \mathcal{O}_\mathcal{X})$.
\begin{enumerate}
\item $H^i_\tau(\mathcal{X}, \mathcal{F}) = 0$ for all $i$.
\item Let $f : \mathcal{X} \to \mathcal{Y}$ be a morphism of algebraic stacks.
Then $R^if_*\mathcal{F}$ (computed in $\tau$-topology) is a
parasitic object of $\textit{Mod}(\mathcal{Y}_\tau, \mathcal{O}_\mathcal{Y})$.
\end{enumerate}
\end{lemma}

\begin{proof}
We first reduce (2) to (1).
By Sheaves on Stacks, Lemma \ref{stacks-sheaves-lemma-pushforward-restriction}
we see that $R^if_*\mathcal{F}$ is the sheaf associated to the presheaf
$$
y \longmapsto
H^i_\tau\Big(V \times_{y, \mathcal{Y}} \mathcal{X},
\ \text{pr}^{-1}\mathcal{F}\Big)
$$
Here $y$ is a typical object of $\mathcal{Y}$ lying over the scheme $V$.
By Lemma \ref{lemma-parasitic} it suffices to show that
these cohomology groups are zero when $y : V \to \mathcal{Y}$ is flat.
Note that $\text{pr} : V \times_{y, \mathcal{Y}} \mathcal{X} \to \mathcal{X}$
is flat as a base change of $y$. Hence by
Lemma \ref{lemma-parasitic} we see that $\text{pr}^{-1}\mathcal{F}$
is parasitic. Thus it suffices to prove (1).

\medskip\noindent
To see (1) we can use the spectral sequence of
Sheaves on Stacks, Proposition
\ref{stacks-sheaves-proposition-smooth-covering-compute-cohomology}
to reduce this to the case where $\mathcal{X}$
is an algebraic stack representable by an algebraic space.
Note that in the spectral sequence each
$f_p^{-1}\mathcal{F} = f_p^*\mathcal{F}$ is a parasitic module by
Lemma \ref{lemma-parasitic} because the morphisms
$f_p : \mathcal{U}_p =
\mathcal{U} \times_\mathcal{X} \ldots
\times_\mathcal{X} \mathcal{U} \to \mathcal{X}$ are flat.
Reusing this spectral sequence one more time (as in the
proof of Lemma \ref{lemma-general-pushforward})
we reduce to the case where the
algebraic stack $\mathcal{X}$ is representable by a scheme $X$.
Then $H^i_\tau(\mathcal{X}, \mathcal{F}) = H^i((\Sch/X)_\tau, \mathcal{F})$.
In this case the vanishing follows easily from an argument
with {\v C}ech coverings, see
Descent, Lemma \ref{descent-lemma-cohomology-parasitic}.
\end{proof}

\end{document}
